\documentclass[11pt]{book}

\usepackage{ppbook}

\begin{document}

\chapter{Introduction}

\section{Why software engineering?}

Software engineering is concerned with the correct development,
implementation and maintenance of software systems.

Some important parts of the software life cycle:

\begin{itemize}
    \item Requirements engineering
    \item Architectural design
    \item Implementation
    \item Testing / verification
    \item Deployment
    \item Long-term maintenance
\end{itemize}

Quality assurance includes formal methods, static analysis,
and unit/integration testing.

Another important idea is abstraction and modular design
$\rightarrow$ divide a system into decoupled components
with well-defined interfaces.

\section{Program analysis}

Given a program $P$, program analysis tries to determine
whether some property holds for all its behaviours.

\[
P \models \phi
\]

For example,

\[
P \models \text{``no divide by zero''}
\]

There are two broad ways of doing analysis:

\begin{itemize}
    \item \textbf{Static} $\rightarrow$ analysis without running
    the program.
    \begin{itemize}
        \item Compiler analysis
        \item Data-flow analysis
        \item Abstract interpretation
        \item Symbolic execution
        \item Model checking
    \end{itemize}

    \item \textbf{Dynamic} $\rightarrow$ observe the program
    while it is running.
    \begin{itemize}
        \item Testing
        \item Profiling
        \item Runtime monitoring
        \item Dynamic taint analysis
    \end{itemize}
\end{itemize}

\section{Formal verification}

Formal verification can be thought of as

\[
\text{Model}+\text{Specification}+\text{Proof}
\]

\begin{ppdefn}
The model describes the semantics of the program.

The specification describes the correctness property
we want to establish.

The proof shows that the program satisfies the specification
under the given semantics.
\end{ppdefn}

Examples include model checking, theorem proving,
and symbolic verification.

\section{Program analysis vs formal verification}

Program analysis may infer useful information about
a program, for example:

\begin{itemize}
    \item amount influences the balance
    \item there are two possible paths
    \item there is no pointer dereference
\end{itemize}

Formal verification goes further and tries to prove
a specified property.

For example,

\[
balance \geq 0 \land amount \geq 0
\Rightarrow
balance' \geq 0
\]

The point is that the property is not just observed
or guessed --- it is proved.

\section{Why this gets difficult}

Think of Rice's theorem.

There is a basic tradeoff:

\[
\text{precision}
\longleftrightarrow
\text{scalability}
\]

More precise analysis generally means more computation
and cost.

Some reasons:

\begin{itemize}
    \item infinitely many possible inputs
    \item number of paths can grow exponentially
    \item loops can create infinitely many executions
\end{itemize}

Real programs also have complicated state involving
pointers, heap objects, data structures, recursion,
callbacks, exceptions, threads, files, networks, etc.

\begin{ppnote}
Exact analysis becomes difficult very quickly when the
program has many possible inputs, paths, and complex
program state.
\end{ppnote}

\section{LLMs and program analysis}

LLMs can help generate possible specifications
from natural language.

\[
\text{Natural language}
\rightarrow
\text{Logic}
\]

For example:

\[
G(req) \rightarrow F(resp)
\]

This can represent the idea that every accepted request
eventually receives a response.

LLMs can also help with semantic inference, such as:

\begin{itemize}
    \item invariants
    \item preconditions
    \item postconditions
    \item auxiliary lemmas
\end{itemize}

Sometimes an invariant is known to the underlying
decision procedure but the SMT solver is unable to
find it by itself.

So an LLM can generate candidate invariants or lemmas.

\begin{ppkey}
LLMs are useful for generating candidate properties,
but those candidates still need to be validated.
\end{ppkey}

\section{Limitations}

LLMs cannot reliably verify their own output,
so their output cannot itself be treated as a proof.

They are neither guaranteed to be sound nor complete.

\[
\text{Soundness}
\rightarrow
\text{no false positives}
\]

\[
\text{Completeness}
\rightarrow
\text{no false negatives}
\]

They can also be weaker at precise multi-step reasoning.

For example,

\[
\forall x\exists y\;R(x,y)
\]

needs the quantifiers to remain in the correct order.
A later reasoning step may accidentally change this.

\section{overall idea}

LLMs can help at several different layers:

\begin{itemize}
    \item \textbf{Property discovery}
    $\rightarrow$ what should be checked?

    \item \textbf{Abstraction discovery}
    $\rightarrow$ what information should the analyzer maintain?

    \item \textbf{Proof search}
    $\rightarrow$ which invariant or lemma makes the proof work?

    \item \textbf{Tool orchestration}
    $\rightarrow$ which analyzer, verifier or solver should be used?
\end{itemize}

\begin{ppkey}
Trust but verify.
\end{ppkey}

\end{document}
